\subsection{Pair partitions of type B}
%\label{subsec:PartitionsB}

Let $S$ be an ordered set. Then $\pi = \{ V_1,\dots, V_p\}$ is a partition of $S$, if the $V_i \neq \varnothing$ are
ordered and disjoint sets $V_i=(v_1,\dots,v_k)$, where $v_1<\dots<v_k$, whose union is $S$.
For $V \in \pi$ , we say that~$V$ is a \emph{block of $\pi$}.
Any partition $\pi$ defines an equivalence relation on $S$,
denoted by $\sim_\pi$, such that the equivalence classes are the
blocks of~$\pi$.
Therefore, $i\sim_\pi j$ if $i$ and $j$ belong to the same block of $\pi$.
A block of $\pi$ is called a \emph{singleton} if it consists of one element.
Similarly, a~block of $\pi$ is called a \emph{pair} if it consists of two elements.
Let $\Sing(\pi)$ and $\Pair(\pi)$ denote the set of all singletons and pairs of $\pi$, respectively.
The set of partitions of $S$ is denoted by $\Part(S)$, in the case where $S =
[n] := \{1, \dots , n\}$ we write $\Part(n)$ := $\Part([n])$.
We denote by $\P_{1,2}(S)$ the subset of partitions $\pi \in \Part(S)$ whose every block is either a
pair or a singleton.

From now on, we will work on a set $[\pm n]:=\{\bar n, \dots , \bar 1, 1,\dots, n\}$.
For a pair $V=(a,b)$ (or a~singleton $V=(a)$), we denote its reflection by $\overline V=(\bar b , \bar a)$ ($\overline V=( \bar a)$), where $\bar a =-a$.
Similarly, we define
\[
\bar \pi: =\big\{\overline{V}\mid V\in \pi,\text{ } \pi \in \P_{1,2}([\pm n])\big\}.
\]

\begin{Definition} \label{def:partycji}
	We denote by $\P_{1,2}^{B}(n)$ the subset of partitions $\pi \in \P_{1,2}([\pm n])$
	such that they are symmetric $\overline{\pi} =
	\pi$ (which is invariant under the bar operation), but every pair $V \in \pi$ is different from its reflection
	$\overline{V}$, i.e.,~$V \neq \overline{V}$.
	We call $\P_{1,2}^{B}(n)$ the set of partitions of type B.
\end{Definition}

From Definition~\ref{def:partycji} it follows that for every block in $B\in \pi$, $\pi \in \P_{1,2}^{B}(n)$ there exists a unique {reflection block} $\bar B\in\pi$.
This leads to one more definition.
We call $\BL= ((\bar b, \bar a),( a,b))$ a \emph{B-pair} if $\bar b <b$, $|{a}| <|b|$ and $(\bar b,\bar a ), (a,b)\in \pi$.
The B-pair $((\bar b, \bar a),( a,b))$ is called \emph{positive} if $b\times a >0 $; otherwise
it is called \emph{negative}; see Figure~\ref{fig:coneced}.
The set of such B-pairs is denoted by $\Pair_B(\pi)$.
Let us notice that from Definition~\ref{def:partycji} it follows that
the element $s$ is a singleton of $\pi \in \P_{1,2}^{B}(n)$ if and only if $\bar s$ is also a singleton of $\pi$, thus we can define the subset of B-singletons $((\bar a),(a))$ as
\[
\Sing_B(\pi):=\{((\bar a), (a)) \mid \bar a,a \in \Sing(\pi), \,\bar a < a\}, \qquad \pi \in \P_{1,2}^{B}(n).
\]
Let
\[
\P_{2}^{B}(n):=\big\{\pi \in \P_{1,2}^{B}(n)\mid \Sing_B(\pi)=\varnothing\big\}
\]
and
\[
\PA_2(n):=\big\{\pi \in \P_{2}^{B}(n)\mid (\bar b,\bar a ),\, (a,b)\in \pi \implies ((\bar b,\bar a ), (a,b))\text{ is positive}\big\},
\]
i.e., it is a subset of $\P_{2}^{B}(n)$, with only positive B-pairs.

\begin{Remark}\quad
\begin{enumerate}
\item [$1.$] We note that a B-pair is not a block of $\pi$, but a pair of reflection pairs.
	
\item [$2.$] We note that $\#\PA_2(n)=(n-1)!!$ for even $n$, which is the same as the number of the classical pair partitions on $n$ points.
	In Figure~\ref{figexample1}, they are depicted on the top row.
\end{enumerate}	
\end{Remark}

\begin{figure}[h]\centering
\MatchingMeandersabc{4}{-4/1, -3/-2,2/3,-1/4 }{}
\caption{The example of $\pi \in \P_{2}^{B}(4)$ with a positive B-pair $((\bar 3,\bar 2 ), (2,3))$ and a negative B-pair $((\bar 4,1), (\bar 1,4))$. }
	\label{fig:coneced}
\end{figure}

We introduce some partition statistics.
For two pairs $V$ and $ W$, we introduce the relation $\text{cr}$ as follows:
\begin{gather*}
V\stackrel{\text{cr}}{\sim}W \iff
V=(i,j),\qquad W=(k,l)\quad \text{such that}\quad i<k<j<l,
\end{gather*}
For a set partition $\pi\in \P_{1,2}^{B}(n)$, let $\Cr(\pi)$ be the number of crossings of B-pairs, i.e.,
\begin{align*}
\Cr(\pi)={}&\#\big\{\big(\big(\overline{V}, V \big),\big(\overline{ W}, W\big)\big)\in \Pair_B(\pi)\times \Pair_B(\pi) \mid \text{$V\stackrel{\text{cr}}{\sim}W $}\big\}
\\
&+\#\big\{\big(\big(\overline{V}, V\big),\big(\overline{ W}, W\big)\big)\in \Pair_B(\pi)\times \Pair_B(\pi) \mid \text{$\overline{V}\stackrel{\text{cr}}{\sim}W $}\big\}.
\end{align*}
For two blocks $V$, $W$ of a set partition, we say that $W$ \emph{covers} $V$ if there are $i,j \in W$ such that $i <k <j$ for any $k\in V$ and then we write $V\stackrel{\text{cs}}{\sim}W $. For $\pi\in\P_{1,2}^{B}(n)$, let $\InS(\pi)$
be the number of pairs of a B-singleton and a covering B-pair:
\begin{align*}
\begin{split}
\InS(\pi)&=\#\big\{\big(\big(\overline{ V}, V \big),\big(\overline{ W}, W\big)
\big) \in \Sing_B(\pi) \times \Pair_B(\pi) \mid V\stackrel{\text{cs}}{\sim}W
\big\}
\\
 & +\#\big\{\big(\big(\overline{V}, V \big),\big(\overline{W}, W\big)\big) \in \Sing_B(\pi) \times \Pair_B(\pi) \mid \overline{V}\stackrel{\text{cs}}{\sim}W\big\}.
 \end{split}
\end{align*}
We say that the B-pair $(\overline{ W} , W)$ cover the B-singleton $\big(\overline{V}, V\big)$ if $V\stackrel{\text{cs}}{\sim}{W} $ or $\overline{V}\stackrel{\text{cs}}{\sim}{ W} $ (equivalently, $\big(\overline{V}, V\big)$ is the inner singleton of $\big(\overline{W} , W\big)$).

Let $\NB(\pi)$ be the number of negative B-pairs of $\Pair_B(\pi)$, where $\pi \in \P_{1,2}^{B}(n)$.
A set partition $\pi \in \P_{1,2}^{B}(n)$ is \emph{non-crossing} if $\Cr(\pi)=0$. The set of non-crossing pair partitions of $[\pm n]$ is denoted by $\NC^B_2(n)$, and by $\NC^A_2(n)$ we denote the subset of $\NC^B_2(n)$ where all B-pairs are positive.

\begin{Remark} \label{jakliczyc}
While reading this, the first impression seems to be that the procedure of counting the crossings is complicated. This is not true since it can be read from the figure as follows. First, we draw a vertical line in the center as in Figure~\ref{fig:examplejakliczyccrosingi}.
Then we count only the crossings on the left of this vertical line.
We count the number of
negative B-pairs as the number of B-pairs crossed by a vertical line ($\scriptstyle{\blacklozenge}$ points on Figure~\ref{fig:examplejakliczyccrosingi}).
Similarly we count the pairs of the form
\[
(\text{B-singleton, covering B-pair})
\]
as the number of such a pairs with at least one leg on the left of this vertical line.

\begin{figure}[htp]
\centering
\MatchingMeanders{10}{-1/3, -3/1 ,-6/5,-5/6,7/10,-10/-7,-2/9,-9/2}
\caption{The example of statistic of partition $\pi \in \P_{1,2}^{B}(10)$, i.e., $\Cr(\pi)=4$, $\NB(\pi)=3$, $\InS(\pi)=5$.}
\label{fig:examplejakliczyccrosingi}
\end{figure}
\end{Remark}

Now we prove the following theorem, which shows the relationship between the set of partitions of type B
(corresponding statistic) and a joint action of Gaussian operators on a vacuum vector.

\begin{Theorem}\label{lem101}
For any $i\in\{1,\ldots,2n\}$ and $x_{\bar i}\otimes x_i \in \HH_\R$, we have
\begin{equation}\label{formula101}
	\state(\G(x_{\overline{ 2n}}\otimes x_{ 2n})\cdots \G(x_{\bar 1}\otimes x_{ 1}))= \sum_{\pi\in\PB_{2}(2n)} \s^{\NB(\pi)}\q^{\Cr(\pi)} \prod_{\substack{(i,j) \in \Pair(\pi)} }\langle x_i, x_j\rangle.
\end{equation}
\end{Theorem}


\begin{proof}
Given
$\epsilon=(\epsilon(1), \dots, \epsilon(n))\in\{1,\ast\}^{n}$, let
$\P_{1,2;\epsilon}^{B}(n)$ be the set of partitions $\pi\in
\P_{1,2}^{B}(n)$ such that
\begin{itemize}\itemsep=0pt
\item if $((\bar b ,\bar a) , (a,b))$ is a B-pair of
$\Pair_B(\pi)$, then $\epsilon(|a|)= \ast$, $\epsilon(|b|)=1$,
\item if $\{c\}$ is a singleton in $\pi$, then $\epsilon( |c|)=\ast$.
\end{itemize}
We will prove first that
\begin{gather}
\B^{\epsilon(n)}(x_{\overline{n}}\otimes x_{n})\cdots \B^{\epsilon(1)}(x_{\overline{1}}\otimes x_1)\Omega \otimes \Omega \nonumber
\\
\qquad{} = \sum_{\pi\in\PB_{1,2;\epsilon}(n)} \s^{\NB(\pi)}q^{\Cr(\pi)+ \InS(\pi)}
\prod_{\substack{(i,j) \in \Pair(\pi)} }\langle x_i, x_j\rangle
\bigotimes_{i\in \Sing(\pi)} x_i\label{formula10112}
\end{gather}
holds, where
\[
\Sing(\pi) = \{(\bar v_m),\dots,(\bar v_1), (v_1),\dots,(v_m)\} \subset\bar{\N} \cup \N, \qquad
\bar v_m<\cdots < \bar v_1< v_1<\cdots < v_m,
\]
and
$\bigotimes_{i\in \Sing(\pi)} x_i$ denotes the tensor product $x_{\bar v_m}\otimes \cdots \otimes x_{\bar v_1} \otimes x_{v_1}\otimes \cdots \otimes x_{v_m}$.
If
\[
\#\{i\in[j] \mid \epsilon(i)=1\}>\#\{i\in[j] \mid \epsilon(i)=\ast\}
\]
for some $j\in\{1,\dots,n\}$, then we understand the sum over the empty set is 0 since $\PB_{1,2;\epsilon}(n)=\varnothing$ in this case.
The proof of \eqref{formula10112} is given by induction.

When $n=1$, then $\B(x_{\bar 1}\otimes x_1)\Omega\otimes \Omega=0$ and $\B^\ast(x_{\bar 1}\otimes x_1)\Omega\otimes \Omega=x_{\bar 1}\otimes x_1$ and hence the formula \eqref{formula10112} is true.

Suppose that the \eqref{formula10112} is true for $n = k$.
We will show that the action of $\B^{\epsilon(k+1)}(x_{\,\overline{k+1}}{}\otimes x_{k+1})$ for $\epsilon(k+1)=\ast$ or $\epsilon(k+1)=1$ corresponds to the inductive pictorial description of set partitions of type B.

\smallskip\noindent
\emph{Case} 1. If $\epsilon(k+1)=\ast$, then the operator $\B^\ast(x_{\,\overline{k+1}}\otimes x_{k+1})$ acts on the tensor product, putting~$x_{k+1}$ on the right and $x_{\,\overline{k+1}}$ on the left. This operation pictorially corresponds to adding the singleton $(\overline{k+1}
)$ and $(k+1)$ to $\pi\in\PB_{1,2;\epsilon}(k)$, to yield the new type-B partition $\tilde{\pi}\in\PB_{1,2;\epsilon}(k+1)$ such that $((\overline{k+1}), (k+1))\in \Sing_B(\tilde{\pi})$. This map $\pi\mapsto \tilde{\pi} $ does not change the numbers $\NB$, $\Cr$ or $\InS$, which is compatible with the fact that the action of $\B^\ast(x_{\,\overline{k+1}}\otimes x_{k+1})$ does not change the coefficient.
Hence the formula~\eqref{formula10112} holds when $n=k+1$ and $\epsilon(k+1)=\ast$.

\smallskip\noindent
\emph{Case} 2. If $\epsilon(k+1)=1$, then we have two subcases, depending on $p_\q$ and $\ell_{\q}$.

If the positive part $p_\q $ (subcase (a)) (resp. $\s \ell_{\q}$ -- subcase (b)) acts on the tensor product on the right hand side of \eqref{formula10112} (for fixed $\pi$), then new $m$ terms appear by using \eqref{rq}. We obtain the equations
\begin{subequations}
\begin{align}
	\begin{split}
		p_\q(x_{\,\overline{k+1}}\otimes x_{k+1}) & x_{\bar v_m}\otimes \cdots \otimes x_{v_m}=
		\sum_{i=1}^m \q^{m-i}\langle x_{\,\overline{k+1}},x_{\bar v_i}\rangle \langle x_{v_i} ,x_{k+1}\rangle\, \\ &\times x_{\bar v_m }\otimes \cdots \otimes \check{x}_{\bar v_i} \otimes \cdots \otimes x_{\bar v_1} \otimes x_{v_1}\otimes \cdots \otimes \check{x}_{v_i} \otimes \cdots \otimes x_m,
	\end{split}\label{eq:equatinpomocdowodgolwny1}
	\\
	\begin{split}
		\alpha\ell_\q(x_{\,\overline{k+1}}\otimes x_{k+1}) & x_{\bar v_m}\otimes \cdots \otimes x_{v_m}=
		\alpha\q^{m-1}\sum_{i=1}^m \q^{i-1}
		\langle x_{\,\overline{k+1}},x_{ v_i}\rangle \langle x_{\bar v_i} ,x_{k+1}\rangle\, \\& \times x_{\bar v_m }\otimes \cdots \otimes \check{x}_{\bar v_i} \otimes \cdots \otimes x_{\bar v_1} \otimes x_{v_1}\otimes \cdots \otimes \check{x}_{v_i} \otimes \cdots \otimes x_m.
	\end{split}\label{eq:equatinpomocdowodgolwny2}
\end{align}
\end{subequations}
Now we will focus on the $i$-th summand of the above equations.

We fix $\pi\in \PB_{1,2;\epsilon}(k)$ and suppose that $\pi$ has singletons
\begin{align*}
\bar v_m<\cdots <\bar v_{ i} < \cdots <\bar v_1 < v_1<\cdots <v_i <\cdots <v_m,\text{ where } i\in [m].
\end{align*}

We assume that $\Pair_B(\pi)$ contains the B-pairs $\big(\overline{ W}_1 , W_1\big),\dots, \big(\overline{ W}_u , W_u\big)$ which cover the B-singleton $((\bar v_i),( v_i))$.
Case~2(a)~-- see Figure~\ref{figtracepositivea}.
In the $i$-th term of \eqref{eq:equatinpomocdowodgolwny1} the inner product $\langle x_{\,\overline{k+1}},x_{\bar v_i}\rangle \langle x_{v_i} ,x_{k+1}\rangle$ appears with coefficient $\q^{m-i}$. Pictorially this corresponds to getting a set partition $\tilde{\pi} \in \PB_{1,2;\epsilon}(k+1)$ by adding the positive B-pair $\big((\overline{k+1}, \overline{v}_i) , (v_i, k+1)\big)$ to $\Pair_B(\tilde \pi)$. This new B-pair crosses the B-pairs $\big(\overline{W}_j , W_j\big)$, $j\in\{1,\dots,u\}$ and so increases the crossing number by~$u$ but decreases the number $\InS(\pi)$ by $u$ because originally $((\bar v_i),( v_i))$ was the inner singleton of B-pairs $\big(\overline{W}_j , W_j\big)$, $j\in\{1,\dots,u\}$.
Now the new inner B-singletons $((\bar v_{i+1}), (v_{i+1})), \dots,(( \bar v_m), (v_m))$ of B-pair $((\overline{k+1}, \overline{v}_i) , (v_i, k+1))$
appear.
Altogether we have
\[
\Cr(\tilde{\pi})=\Cr(\pi)+u, \qquad
\InS(\tilde{\pi})=\InS(\pi)+m-i-u\qquad \text{and}\qquad
\NB(\tilde{\pi})=\NB(\pi).
\]
So the exponent of $\q$ increases by $m-i$. This factor $\q^{m-i}$ is exactly the factor appearing in equation~\eqref{eq:equatinpomocdowodgolwny1}.

\begin{figure}[h]
\centering
\MatchingMeandersPositiveTracea{9}{-9/-8,8/9,-6/2,-2/6 }{}
\MatchingMeandersPositiveTraceb{10}{-9/-8,8/9,-6/2,-2/6 ,-10/-4,4/10 }{}
\caption{The visualization of the action of $p_\q(x_{\,\overline{k+1}}\otimes x_{k+1})$.}
\label{figtracepositivea}
\end{figure}

Case 2(b) -- see Figure~\ref{figtracepositive}.
In the $i$-th term of \eqref{eq:equatinpomocdowodgolwny2} the inner product $\langle x_{\,\overline{k+1}},x_{ v_i}\rangle \langle x_{\bar v_i} ,x_{k+1}\rangle$ appears with the coefficient $\s \q^{m+ i-2}$.
Graphically this
corresponds to getting a set partition $\tilde{\pi} \in \PB_{1,2;\epsilon}(k+1)$ by adding $\overline{k + 1}$ and $k + 1$ to $\pi$ and creating the new negative B-pair $\big((\overline{k+1}, v_{i}) , (\overline v_ i, k+1)\big)\in \Pair_B(\tilde \pi)$.
Similarly to Case 2(a), we count the change of numbers and get
\[
\Cr(\tilde{\pi})=\Cr(\pi)+u,\quad
\InS(\tilde{\pi})=\InS(\pi)+m-1+i-1-u\quad\text{and}\quad \NB(\tilde{\pi})=\NB(\pi)+1.
\]
Altogether, when moving from $\pi$ to $\tilde{\pi}$, the exponent of $\s$ increases by $1$ and the exponent of $\q$ increases by ${m+ i-2}$, which coincides with the coefficient appearing in the action of $\s \ell_{\q}(x_{\,\overline{k+1}}\otimes x_{k+1})$, creating the inner product $\langle x_{\,\overline{k+1}},x_{ v_i}\rangle \langle x_{\bar v_i}, x_{k+1}\rangle$.

\begin{figure}[h]
%\begin{center}
%\QQ{$(b)$}
%\end{center}%
\centering
\MatchingMeandersPositiveTracea{9}{-9/-8,8/9,-6/2,-2/6 }{} % $\xrightarrow{\text{trace action }}$
\MatchingMeandersPositiveTraceb{10}{-9/-8,8/9,-6/2,-2/6 ,-10/4,-4/10 }{}
\caption{The visualization of the action of $\s \ell_{\q}(x_{\,\overline{k+1}}\otimes x_{k+1})$.}
\label{figtracepositive}
\end{figure}


Note that as $\pi$ runs over $\PB_{1,2;(\epsilon(1),\dots,\epsilon(k))}(k)$, every set partition $\tilde{\pi}\in \PB_{1,2;(\epsilon(1),\dots,\epsilon(k),1)}(k+1)$ appears exactly once, either in Case~2(a) or in Case~2(b). Therefore, in Case~2, the pictorial inductive step and the actual action of $\B(x_{\,\overline{k+1}}\otimes x_{k+1})$ both create the same terms with the same coefficients, and hence the formula \eqref{formula10112} is true when $n=k+1$ and $\epsilon(k+1)=1$. Case~1 and Case~2 show by induction that the formula \eqref{formula10112} holds for all $n \in \N$.

For a given set of B-pairs $\{((\bar b_1, \bar a_1),(a_1,b_1)), \dots, ((\bar b_{n/2}, \bar a_{n/2}),(a_{n/2},b_{n/2}))\}\in \Pair_B( \pi) $,
where $\pi \in \PB_{2}(n)$, we denote the set of \emph{left and right legs} of $\pi$ by $l_\pi :=\{ a_1,\bar a_1,\dots, a_{n/2}, \bar a_{n/2} \}$, and $r_\pi :=\{ b_1,\bar b_1,\dots, b_{n/2}, \bar b_{n/2} \}$. Formula \eqref{formula101} follows from \eqref{formula10112} by taking the sum over all $\epsilon$ such that $\Sing(\pi)=\varnothing$. In this case, we understand that $\bigotimes_{i\in \Sing(\pi)} x_i=\Omega\otimes \Omega $ and
\begin{align*}
\bigsqcup_{\e\in\{1,\ast\}^{n}} \PB_{2;\e}(n)&=\bigsqcup_{\e\in\{1,\ast\}^{n}} \big\{\pi \in \P_{2}^{B}(n)\mid (\bar b,\bar a ), (a,b)\in \pi \implies \epsilon(|a|)= \ast,\, \epsilon(|b|)=1\big\}
\\
& =\bigsqcup_{\substack{ L\subset [\pm n] \\ \# L =n
}} \big\{\pi \in \P_{2}^{B}(n)\mid
l_\pi=L \text{ and }r_\pi = [\pm n]\setminus L\big\}
\\
&=\PB_{2}(n).
\end{align*}

 Finally, applying the state action, we get
\begin{align*}
&\state(\G(x_{\overline{ 2n}}\otimes x_{ 2n})\cdots \G(x_{\bar 1}\otimes x_{ 1}))
\\
&\qquad =\sum_{\e\in\{1,\ast\}^{2n}} \state\big( \B^{\epsilon(2n)}(x_{\overline{2n}}\otimes x_{2n})\cdots \B^{\epsilon(1)}(x_{\overline{1}}\otimes x_1)\big)
\\
&\qquad =\sum_{\e\in\{1,\ast\}^{2n}} \sum_{\substack{\pi\in\PB_{1,2;\epsilon}(2n)\\ \Sing(\pi)=\varnothing }} \s^{\NB(\pi)}q^{\Cr(\pi)+ \InS(\pi)}
\prod_{\substack{(i,j) \in \Pair(\pi)} }\langle x_i, x_j\rangle
\\
&\qquad =\sum_{\e\in\{1,\ast\}^{2n}} \sum_{\substack{\pi\in\PB_{2;\epsilon}(2n) }} \s^{\NB(\pi)}q^{\Cr(\pi)}
\prod_{\substack{(i,j) \in \Pair(\pi)} }\langle x_i, x_j\rangle
\\
&\qquad = \sum_{\pi\in\PB_{2}(2n)} \s^{\NB(\pi)}\q^{\Cr(\pi)} \prod_{\substack{(i,j) \in \Pair(\pi)} }\langle x_i, x_j\rangle.
\tag*{\qed}
\end{align*}
\renewcommand{\qed}{}
\end{proof}
